% Options for packages loaded elsewhere
% Options for packages loaded elsewhere
\PassOptionsToPackage{unicode}{hyperref}
\PassOptionsToPackage{hyphens}{url}
\PassOptionsToPackage{dvipsnames,svgnames,x11names}{xcolor}
%
\documentclass[
  japanese,
  lualatex,
  ja=standard,
  12pt,
  a4paper]{ltjsarticle}
\usepackage{xcolor}
\usepackage{amsmath,amssymb}
\setcounter{secnumdepth}{5}
\usepackage{iftex}
\ifPDFTeX
  \usepackage[T1]{fontenc}
  \usepackage[utf8]{inputenc}
  \usepackage{textcomp} % provide euro and other symbols
\else % if luatex or xetex
  \usepackage{unicode-math} % this also loads fontspec
  \defaultfontfeatures{Scale=MatchLowercase}
  \defaultfontfeatures[\rmfamily]{Ligatures=TeX,Scale=1}
\fi
\usepackage{lmodern}
\ifPDFTeX\else
  % xetex/luatex font selection
\fi
% Use upquote if available, for straight quotes in verbatim environments
\IfFileExists{upquote.sty}{\usepackage{upquote}}{}
\IfFileExists{microtype.sty}{% use microtype if available
  \usepackage[]{microtype}
  \UseMicrotypeSet[protrusion]{basicmath} % disable protrusion for tt fonts
}{}
\makeatletter
\@ifundefined{KOMAClassName}{% if non-KOMA class
  \IfFileExists{parskip.sty}{%
    \usepackage{parskip}
  }{% else
    \setlength{\parindent}{0pt}
    \setlength{\parskip}{6pt plus 2pt minus 1pt}}
}{% if KOMA class
  \KOMAoptions{parskip=half}}
\makeatother
% Make \paragraph and \subparagraph free-standing
\makeatletter
\ifx\paragraph\undefined\else
  \let\oldparagraph\paragraph
  \renewcommand{\paragraph}{
    \@ifstar
      \xxxParagraphStar
      \xxxParagraphNoStar
  }
  \newcommand{\xxxParagraphStar}[1]{\oldparagraph*{#1}\mbox{}}
  \newcommand{\xxxParagraphNoStar}[1]{\oldparagraph{#1}\mbox{}}
\fi
\ifx\subparagraph\undefined\else
  \let\oldsubparagraph\subparagraph
  \renewcommand{\subparagraph}{
    \@ifstar
      \xxxSubParagraphStar
      \xxxSubParagraphNoStar
  }
  \newcommand{\xxxSubParagraphStar}[1]{\oldsubparagraph*{#1}\mbox{}}
  \newcommand{\xxxSubParagraphNoStar}[1]{\oldsubparagraph{#1}\mbox{}}
\fi
\makeatother


\usepackage{longtable,booktabs,array}
\newcounter{none} % for unnumbered tables
\usepackage{calc} % for calculating minipage widths
% Correct order of tables after \paragraph or \subparagraph
\usepackage{etoolbox}
\makeatletter
\patchcmd\longtable{\par}{\if@noskipsec\mbox{}\fi\par}{}{}
\makeatother
% Allow footnotes in longtable head/foot
\IfFileExists{footnotehyper.sty}{\usepackage{footnotehyper}}{\usepackage{footnote}}
\makesavenoteenv{longtable}
\usepackage{graphicx}
\makeatletter
\newsavebox\pandoc@box
\newcommand*\pandocbounded[1]{% scales image to fit in text height/width
  \sbox\pandoc@box{#1}%
  \Gscale@div\@tempa{\textheight}{\dimexpr\ht\pandoc@box+\dp\pandoc@box\relax}%
  \Gscale@div\@tempb{\linewidth}{\wd\pandoc@box}%
  \ifdim\@tempb\p@<\@tempa\p@\let\@tempa\@tempb\fi% select the smaller of both
  \ifdim\@tempa\p@<\p@\scalebox{\@tempa}{\usebox\pandoc@box}%
  \else\usebox{\pandoc@box}%
  \fi%
}
% Set default figure placement to htbp
\def\fps@figure{htbp}
\makeatother



\ifLuaTeX
\usepackage[bidi=basic,shorthands=off,provide=*]{babel}
\else
\usepackage[bidi=default,shorthands=off,provide=*]{babel}
\fi
\ifLuaTeX
  \usepackage{selnolig} % disable illegal ligatures
\fi


\setlength{\emergencystretch}{3em} % prevent overfull lines

\providecommand{\tightlist}{%
  \setlength{\itemsep}{0pt}\setlength{\parskip}{0pt}}



 


\makeatletter
\@ifpackageloaded{tcolorbox}{}{\usepackage[skins,breakable]{tcolorbox}}
\@ifpackageloaded{fontawesome5}{}{\usepackage{fontawesome5}}
\definecolor{quarto-callout-color}{HTML}{909090}
\definecolor{quarto-callout-note-color}{HTML}{0758E5}
\definecolor{quarto-callout-important-color}{HTML}{CC1914}
\definecolor{quarto-callout-warning-color}{HTML}{EB9113}
\definecolor{quarto-callout-tip-color}{HTML}{00A047}
\definecolor{quarto-callout-caution-color}{HTML}{FC5300}
\definecolor{quarto-callout-color-frame}{HTML}{acacac}
\definecolor{quarto-callout-note-color-frame}{HTML}{4582ec}
\definecolor{quarto-callout-important-color-frame}{HTML}{d9534f}
\definecolor{quarto-callout-warning-color-frame}{HTML}{f0ad4e}
\definecolor{quarto-callout-tip-color-frame}{HTML}{02b875}
\definecolor{quarto-callout-caution-color-frame}{HTML}{fd7e14}
\makeatother
\makeatletter
\@ifpackageloaded{caption}{}{\usepackage{caption}}
\AtBeginDocument{%
\ifdefined\contentsname
  \renewcommand*\contentsname{目次}
\else
  \newcommand\contentsname{目次}
\fi
\ifdefined\listfigurename
  \renewcommand*\listfigurename{図一覧}
\else
  \newcommand\listfigurename{図一覧}
\fi
\ifdefined\listtablename
  \renewcommand*\listtablename{表一覧}
\else
  \newcommand\listtablename{表一覧}
\fi
\ifdefined\figurename
  \renewcommand*\figurename{図}
\else
  \newcommand\figurename{図}
\fi
\ifdefined\tablename
  \renewcommand*\tablename{表}
\else
  \newcommand\tablename{表}
\fi
}
\@ifpackageloaded{float}{}{\usepackage{float}}
\floatstyle{ruled}
\@ifundefined{c@chapter}{\newfloat{codelisting}{h}{lop}}{\newfloat{codelisting}{h}{lop}[chapter]}
\floatname{codelisting}{コード}
\newcommand*\listoflistings{\listof{codelisting}{コード一覧}}
\makeatother
\makeatletter
\makeatother
\makeatletter
\@ifpackageloaded{caption}{}{\usepackage{caption}}
\@ifpackageloaded{subcaption}{}{\usepackage{subcaption}}
\makeatother
\usepackage{bookmark}
\IfFileExists{xurl.sty}{\usepackage{xurl}}{} % add URL line breaks if available
\urlstyle{same}
% fallback for those not using the hyperref driver hyperxmp:
\makeatletter
\@ifundefined{xmpquote}{\newcommand{\xmpquote}[1]{#1}}{}
\makeatother
\hypersetup{
  pdftitle={第4回：資本蓄積のないRBCモデル},
  pdfauthor={宮﨑憲治},
  pdflang={ja},
  colorlinks=true,
  linkcolor={blue},
  filecolor={Maroon},
  citecolor={Blue},
  urlcolor={Blue},
  pdfcreator={LaTeX via pandoc}}


\title{第4回：資本蓄積のないRBCモデル}
\usepackage{etoolbox}
\makeatletter
\providecommand{\subtitle}[1]{% add subtitle to \maketitle
  \apptocmd{\@title}{\par {\large #1 \par}}{}{}
}
\makeatother
\subtitle{マクロ経済学B}
\author{宮﨑憲治}
\date{2026年度}
\begin{document}
\maketitle

\renewcommand*\contentsname{目次}
{
\setcounter{tocdepth}{3}
\tableofcontents
}

\section{講義の目的}\label{ux8b1bux7fa9ux306eux76eeux7684}

\subsection{本講義で探究する中心的な問い}\label{ux672cux8b1bux7fa9ux3067ux63a2ux7a76ux3059ux308bux4e2dux5fc3ux7684ux306aux554fux3044}

本講義では、経済全体の資本ストック \(\bar{K}\) を \(1\)
に固定（正規化）し、投資による資本蓄積を行わない最も基本的な「\textbf{柔軟価格（Flexible
Price）モデル}」を構築・分析します。標準的なRBC（Real Business
Cycle）モデルにおいて重要な役割を果たす投資の通時的調整が存在しないため、経済のダイナミクスは極めてシンプルかつ透明になります。すなわち、外生的な技術ショック（生産性ショック）が観察された後、当期の家計による労働供給と企業による労働需要の相互作用によって、当期の実物配分（産出、消費、労働、実質賃金）が瞬時に決定されます。

この経済の静学的配分を特徴づける中心的な均衡条件は、以下の2式に集約されます。
\[
MRS_t = W_t = MPN_t, \qquad Y_t = C_t.
\] -
第1式は、家計の主観的な\textbf{限界代替率（MRS）}、市場の\textbf{実質賃金（\(W_t\)）}、および企業の技術的な\textbf{労働の限界生産力（MPN）}が完全に一致するという労働市場の清算条件です。
-
第2式は、投資も政府支出も存在しない閉鎖経済における財市場の\textbf{需給均衡（資源制約式）}です。

標準的な効用関数（CRRA効用・Frisch労働不効用）とCobb-Douglas生産関数を特定化すると、外生的な技術ショック
\(a_t\)
に対する各マクロ変数の対数線形反応は、以下の明快な解析解として導かれます。
\[
\begin{aligned}
y_t &= \zeta_a a_t,\\
n_t &= \frac{1-\gamma}{\varphi+\alpha+\gamma(1-\alpha)} a_t,\\
r_t &= \gamma \zeta_a (\rho_a - 1) a_t.
\end{aligned}
\] ここで、技術進歩に対して労働投入 \(n_t\)
が増加するか減少するかは、「賃金上昇に伴い働くインセンティブが高まる代替効果」と「富が増えて余暇を増やそうとする所得効果」の相対的な大小関係（パラメータ
\(\gamma \gtrless 1\)）によって鮮やかに決まります。また、実質利子率
\(r_t\)
は、実物側で先に決まった消費経路を純供給ゼロの金融市場と整合させる「\textbf{自然実質利子率（Natural
Rate of Interest）}」として受動的に復元されます。

本講義の到達目標は、以下の4点です。

\begin{enumerate}
\def\labelenumi{\arabic{enumi}.}
\tightlist
\item
  \textbf{静学的均衡条件の導出}：\(MRS_t = W_t = MPN_t\)
  および財市場清算から、経済の実物配分を一意に決定する論理構造を理解する。
\item
  \textbf{ショック反応の解析}：技術ショックに対する産出、消費、労働、実質賃金の反応係数を導出し、代替効果と所得効果のメカニズムを説明できるようになる。
\item
  \textbf{自然実質利子率の復元}：実物側で決定された消費成長経路から、消費オイラー方程式を通じて自然実質利子率を導出できる。
\item
  \textbf{モデル拡張への足がかり}：第5回で政府支出を導入した際、どの均衡条件がどのように修正されるかを明確に予見できるようにする。
\end{enumerate}

図~\ref{fig-lecture04-overview}
は、家計と企業の静学的最適化条件から実物配分が決定され、それが自然実質利子率の導出へとつながる理論的展開を視覚的に整理したものです。

\begin{figure}[H]

\centering{

\includegraphics[width=0.95\linewidth,height=\textheight,keepaspectratio]{lecture04_files/mediabag/figures/lecture04/lecture04_overview_infographic.pdf}

}

\caption{\label{fig-lecture04-overview}第4回講義の全体像：固定資本のもとで、家計の労働供給と企業の労働需要が当期の実物配分を一意に決定し、オイラー方程式から自然実質利子率が導出される論理構造。}

\end{figure}%

\section{表記と基本設定}\label{ux8868ux8a18ux3068ux57faux672cux8a2dux5b9a}

本講義ノートでは、慣例に従い、変数の\textbf{水準値（Level）を大文字}で、定常状態からの\textbf{対数偏差（Log-deviation）を小文字}で表記します。たとえば、産出水準
\(Y_t\) とその定常状態値 \(Y\) に対し、対数偏差は
\(y_t \equiv \ln(Y_t / Y) \approx (Y_t - Y)/Y\)
と定義されます。同様に、\(c_t, n_t, w_t, r_t\)
もそれぞれ消費、労働、実質賃金、総実質利子率の対数偏差を表します。

本ノートで使用する主要な変数の記号一覧は以下の通りです。

{\def\LTcaptype{none} % do not increment counter
\begin{longtable}[]{@{}
  >{\raggedright\arraybackslash}p{(\linewidth - 6\tabcolsep) * \real{0.2500}}
  >{\raggedright\arraybackslash}p{(\linewidth - 6\tabcolsep) * \real{0.2500}}
  >{\raggedright\arraybackslash}p{(\linewidth - 6\tabcolsep) * \real{0.2500}}
  >{\raggedright\arraybackslash}p{(\linewidth - 6\tabcolsep) * \real{0.2500}}@{}}
\toprule\noalign{}
\begin{minipage}[b]{\linewidth}\raggedright
記号
\end{minipage} & \begin{minipage}[b]{\linewidth}\raggedright
経済学的意味
\end{minipage} & \begin{minipage}[b]{\linewidth}\raggedright
記号
\end{minipage} & \begin{minipage}[b]{\linewidth}\raggedright
経済学的意味
\end{minipage} \\
\midrule\noalign{}
\endhead
\bottomrule\noalign{}
\endlastfoot
\(C_t\) & 実質消費水準 & \(N_t\) & 労働投入量（就業時間） \\
\(Y_t\) & 総産出水準（実質GDP） & \(W_t\) & 実質賃金率 \\
\(\bar{K}=1\) & 固定された実物資本ストック（\(1\) に正規化） & \(a_t\) &
技術水準（全要素生産性ショック）の対数値 \\
\(B_t\) & 民間実質債券の保有残高 & \(R_t\) &
1期間実質総利子率（グロス利子率） \\
\(\Theta_t\) & 企業の株式（持分）保有株数 & \(V_t\) &
配当落ち後の企業株式時価総額（株価） \\
\(D_t\) & 企業が支払う実質配当（企業利潤） & \(R_t^N\) &
1期間名目総利子率 \\
\(\Pi_t\) & 粗インフレ率（\(P_t/P_{t-1}\)） & \(\mathbb{E}_t\) & 時点
\(t\) の情報集合に基づく条件付き期待値演算子 \\
\end{longtable}
}

また、経済の構造を規定する基礎パラメータおよび頻出する派生パラメータは以下の通りです。

{\def\LTcaptype{none} % do not increment counter
\begin{longtable}[]{@{}
  >{\raggedright\arraybackslash}p{(\linewidth - 2\tabcolsep) * \real{0.5000}}
  >{\raggedright\arraybackslash}p{(\linewidth - 2\tabcolsep) * \real{0.5000}}@{}}
\toprule\noalign{}
\begin{minipage}[b]{\linewidth}\raggedright
記号
\end{minipage} & \begin{minipage}[b]{\linewidth}\raggedright
定義・経済学的意味
\end{minipage} \\
\midrule\noalign{}
\endhead
\bottomrule\noalign{}
\endlastfoot
\(0 < \beta < 1\) & 家計の主観的割引因子 \\
\(\gamma > 0\) & 相対的リスク回避度（異時点間消費代替弾力性の逆数
\(1/\gamma\)） \\
\(\varphi > 0\) & 労働の限界不効用パラメータ（Frisch労働供給弾力性の逆数
\(1/\varphi\)） \\
\(0 \leq \alpha < 1\) &
生産における固定資本の所得分配率（資本シェア） \\
\(0 \leq \rho_a < 1\) & 技術ショックの自己相関（持続性）パラメータ \\
\(0 < \bar{L} < \infty\) &
純金融資産に対する所与の借入限度枠（信用制約下限 \(-\bar{L}\)） \\
\(MRS_t \equiv -U_{N,t}/U_{C,t}\) &
主観的限界代替率（消費財で評価した労働の限界不効用） \\
\(MPN_t \equiv \exp(a_t) F_N(N_t)\) & 労働の限界生産力 \\
\(Q_{t,t+1} \equiv \beta U_{C,t+1}/U_{C,t}\) &
家計の限界効用に基づく確率的割引因子（SDF） \\
\(\zeta_a \equiv \dfrac{1+\varphi}{\varphi+\alpha+\gamma(1-\alpha)}\) &
産出の技術ショックに対する反応弾力性係数 \\
\end{longtable}
}

なお、経済全体における家計および企業の総測度（人口）、ならびに発行株式総数はすべて
\(1\) に正規化されています。したがって、1株あたり株価 \(V_t\)
は企業全体の株式時価総額と完全に一致します。また、固定資本は企業が物理的に保有しており、家計は企業の株式を保有することを通じて、資本の生み出す収益を配当として間接的に受け取ります。

\section{均衡条件}\label{ux5747ux8861ux6761ux4ef6}

\subsection{家計の最適化行動}\label{ux5bb6ux8a08ux306eux6700ux9069ux5316ux884cux52d5}

各期 \(t\) の期首において、まず外生的な技術ショック \(a_t\)
が実現・観察されます。家計はその情報を織り込んだ上で、当期の消費
\(C_t\)、労働供給 \(N_t\)、翌期への無リスク債券保有額
\(B_t\)、および株式保有量 \(\Theta_t\)
を決定します。期首時点で前期から引き継ぐ資産残高 \(B_{t-1}(h)\) および
\(\Theta_{t-1}(h)\) は所与の外生値です。

すべての家計が同一の選好を持ち、初期保有資産が
\(B_{-1}(h)=0, \Theta_{-1}(h)=1\)
で等しい「\textbf{代表的家計（対称均衡）}」を想定します。代表的家計
\(h\) の解くべき動学的効用最大化問題は、以下のように定式化されます。 \[
\max_{\{C_t(h), N_t(h), B_t(h), \Theta_t(h)\}_{t=0}^{\infty}} \mathbb{E}_0 \sum_{t=0}^{\infty} \beta^t U(C_t(h), N_t(h))
\] 各期における実質フロー予算制約式は次式で与えられます。 \[
C_t(h) + B_t(h) + \Theta_t(h) V_t = R_{t-1} B_{t-1}(h) + \Theta_{t-1}(h) (V_t + D_t) + W_t N_t(h).
\]
左辺は家計の当期総支出（消費、新規の債券購入、新規の株式購入）であり、右辺は利用可能な総実質資源（前期債券の元利合計、前期から持ち越した株式の売却評価額および受取配当、当期の労働所得）です。ここで
\(R_{t-1}\) は前期に決定された実質総利子率です。また \(V_t\)
は配当落ち後の株価であるため、前期からの株式保有者は当期の配当 \(D_t\)
と当期の株価 \(V_t\) の双方を享受します。

ポンジ・スキームを排除するため、家計の期末純金融資産に対して以下の信用限度（借入上限）を設けます。
\begin{equation}\protect\phantomsection\label{eq-financial-credit-limit}{
A_t^F(h) \equiv B_t(h) + V_t \Theta_t(h) \geq -\bar{L}, \qquad 0 < \bar{L} < \infty.
}\end{equation}
本講義では、この信用制約が厳密にはバインドせず、消費と労働の選択が常に実行可能領域の内点にある標準的な対称均衡を考察します。

第3回で学んだ通り、家計の異時点間限界代替率から定まる確率的割引因子（SDF）を
\[
Q_{t,t+1} \equiv \beta \frac{U_{C,t+1}}{U_{C,t}}
\]
と定義すると、家計の最適化の一階条件（FOC）は以下の3式として得られます。
\begin{equation}\protect\phantomsection\label{eq-household-focs}{
\begin{aligned}
W_t &= -\frac{U_{N,t}}{U_{C,t}},\\
1 &= \mathbb{E}_t \left[ Q_{t,t+1} R_t \right],\\
V_t &= \mathbb{E}_t \left[ Q_{t,t+1} (V_{t+1} + D_{t+1}) \right].
\end{aligned}
}\end{equation} -
第1式は\textbf{静学的な労働供給条件}であり、労働を1単位増やすことの主観的限界不効用（消費単位換算：\(-U_{N,t}/U_{C,t}\)）と、労働の対価として得られる実質賃金
\(W_t\) を一致させることを意味します。 -
第2式は\textbf{無リスク実質債券に関する消費オイラー方程式}であり、当期の限界効用と翌期の割引期待限界効用（利子率勘案後）を結びつけます。
- 第3式は\textbf{株式の資産価格付け式}であり、当期の株価 \(V_t\)
が、翌期の配当と配当落ち株価の和（\(V_{t+1}+D_{t+1}\)）を確率的割引因子
\(Q_{t,t+1}\) で割り引いた期待現在価値と一致することを示しています。

さらに、第3回と同様に、期末の純金融資産に対して以下の横断性条件（TVC）を課します。
\begin{equation}\protect\phantomsection\label{eq-household-tvc}{
\lim_{k \to \infty} \mathbb{E}_t \left[ Q_{t,t+k} \left( B_{t+k} + V_{t+k} \Theta_{t+k} \right) \right] = 0.
}\end{equation} ここで \(Q_{t,t+k} \equiv \beta^k U_{C,t+k}/U_{C,t}\) は
\(k\)
期間にわたる確率的割引因子です。この条件は、無限の将来において家計の保有資産価値が発散しないことを保証し、合理的な資産価格形成の前提となります。

\begin{tcolorbox}[enhanced jigsaw, arc=.35mm, bottomrule=.15mm, bottomtitle=1mm, breakable, colback=white, colbacktitle=quarto-callout-note-color!10!white, colframe=quarto-callout-note-color-frame, coltitle=black, left=2mm, leftrule=.75mm, opacityback=0, opacitybacktitle=0.6, rightrule=.15mm, title=\textcolor{quarto-callout-note-color}{\faInfo}\hspace{0.5em}{第3回との接続：株式価値}, titlerule=0mm, toprule=.15mm, toptitle=1mm]

第3回の配当割引モデルと同様に、投機的バブルが存在しない「無バブル条件（Bubble-free
Condition）」のもとでは、一階条件の第3式を前向きに逐次代入することにより、株式価値
\(V_t\)
は将来にわたる実質配当流列の割引現在価値の総和として一意に表現されます。
\begin{equation}\protect\phantomsection\label{eq-fundamental-equity-value}{
V_t = \sum_{k=1}^{\infty} \mathbb{E}_t [Q_{t,t+k} D_{t+k}].
}\end{equation}
本講義では、このファンダメンタルズに基づく無バブル価格を株価として採用します。

\end{tcolorbox}

\subsection{企業の利潤最大化行動}\label{ux4f01ux696dux306eux5229ux6f64ux6700ux5927ux5316ux884cux52d5}

代表的な完全競争企業は、所与の固定資本ストック \(\bar{K} = 1\)
と技術水準 \(a_t\) のもとで、毎期の利潤を最大化するように労働投入量
\(N_t\)
を決定します。本モデルでは投資による資本蓄積が行われず、資本レンタル市場も存在しないため、企業の意思決定には将来へ持ち越される状態変数が存在しません。したがって、企業の価値最大化問題は各期ごとの静学的な利潤最大化問題へと自然に帰着されます。

企業の生産関数および当期の実質利潤（実質配当）\(D_t\)
は、次式で定義されます。 \[
Y_t = \exp(a_t) \bar{K}^\alpha N_t^{1-\alpha} = \exp(a_t) N_t^{1-\alpha}, \qquad D_t = Y_t - W_t N_t.
\] 完全競争のもとでプライス・テイカーとして行動する企業は、当期利潤
\(D_t = \exp(a_t) N_t^{1-\alpha} - W_t N_t\) を労働投入量 \(N_t\)
に関して最大化します。その一階条件は、労働の限界生産力（\(MPN_t\)）を市場の実質賃金率
\(W_t\) に等置することです。
\begin{equation}\protect\phantomsection\label{eq-firm-conditions}{
\begin{aligned}
W_t &= (1-\alpha) \exp(a_t) N_t^{-\alpha} = MPN_t,\\
D_t &= Y_t - W_t N_t = \alpha Y_t.
\end{aligned}
}\end{equation} 労働需要条件 \(W_t = (1-\alpha) Y_t / N_t\)
より、企業から労働者への総賃金支払い額は \(W_t N_t = (1-\alpha) Y_t\)
となります。したがって、生産された総産出 \(Y_t\) のうち、\((1-\alpha)\)
の割合は労働への報酬（労働所得）として支払われ、残る \(\alpha\)
の割合は固定資本の所有者である株主への実質配当（資本所得）として分配されます。

\subsection{集計均衡と財市場清算}\label{ux96c6ux8a08ux5747ux8861ux3068ux8ca1ux5e02ux5834ux6e05ux7b97}

民間債券は家計間でのみ貸借される純供給ゼロの資産であり、企業の株式発行総数は
\(1\)
に正規化されています。したがって、金融市場の清算条件は以下の通りです。
\[
\int_0^1 B_t(h)\,dh = 0, \qquad \int_0^1 \Theta_t(h)\,dh = 1.
\]
各家計のフロー予算制約式を経済全体で集計すると、純供給ゼロの債券項は相殺されて消滅し、株式の売買代金も相殺されます。その結果、集計された家計の総支出と総所得の関係は
\[
C_t = W_t N_t + D_t
\] となります。ここに企業利潤の定義式 \(D_t = Y_t - W_t N_t\)
を代入すると、投資も政府支出も存在しない閉鎖経済における財市場の需給均衡条件（資源制約式）
\[
Y_t = C_t
\] が直ちに導かれます。

対称均衡においては、すべての家計が同一の選択を行うため、各家計の保有量も
\(B_t = 0, \Theta_t = 1\) となります。企業の配当
\(D_t = \alpha Y_t \geq 0\)
は常に非負であり、それを現在価値に割り引いた株価 \(V_t\)
も非負であるため、家計の純金融資産は \(A_t^F = V_t \geq 0 > -\bar{L}\)
となり、信用制約は厳密に内点で満たされます。

以上より、資本蓄積および政府部門が存在しない本モデルにおいて、各期の実物配分（\(W_t, N_t, Y_t, C_t\)）を決定する静学的均衡体系は、以下の4本の方程式として完結します。
\begin{equation}\protect\phantomsection\label{eq-static-equilibrium}{
\begin{aligned}
W_t &= -\frac{U_{N,t}}{U_{C,t}},\\
W_t &= (1-\alpha) \exp(a_t) N_t^{-\alpha},\\
Y_t &= \exp(a_t) N_t^{1-\alpha},\\
Y_t &= C_t.
\end{aligned}
}\end{equation}

家計の労働供給条件と企業の労働需要条件を等置することにより、実物配分の核心をなす\textbf{労働市場の均衡条件}が得られます。
\begin{equation}\protect\phantomsection\label{eq-labor-market}{
MRS_t = W_t = MPN_t.
}\end{equation}
この中心的な関係式は、家計が追加的に1単位働くために要求する主観的な消費補償（限界代替率
\(MRS_t\)）と、その追加的労働が技術的に生み出す生産物（限界生産力
\(MPN_t\)）が、市場の実質賃金率 \(W_t\)
を媒介として完全に一致することを表しています。技術ショック \(a_t\)
が与えられると、まずこの労働市場の均衡から当期の労働投入量 \(N_t\)
が決定され、次いで生産関数と資源制約式から産出量 \(Y_t\) および消費水準
\(C_t\) が順次決定されます。

\section{関数形と閉形式解}\label{ux95a2ux6570ux5f62ux3068ux9589ux5f62ux5f0fux89e3}

\subsection{効用関数と生産技術の特定化}\label{ux52b9ux7528ux95a2ux6570ux3068ux751fux7523ux6280ux8853ux306eux7279ux5b9aux5316}

モデルの分析を具体的に進めるため、マクロ経済学において最も標準的に用いられる関数形を導入します。

家計の消費に関する期間効用関数 \(u(C)\)
として、\textbf{一定相対的リスク回避度（CRRA: Constant Relative Risk
Aversion）効用関数}を採用します。 \[
u(C) =
\begin{cases}
\dfrac{C^{1-\gamma}-1}{1-\gamma}, & \gamma \neq 1,\\[8pt]
\ln C, & \gamma = 1.
\end{cases}
\] 各期の期間効用関数 \(U(C_t, N_t)\)
は、消費による効用と労働による不効用が加法分離された次式で与えられます。
\[
U(C_t, N_t) = u(C_t) - \frac{N_t^{1+\varphi}}{1+\varphi}, \qquad \gamma > 0, \quad \varphi > 0.
\] 企業の生産技術は、固定資本 \(\bar{K}=1\)
のもとでの\textbf{Cobb-Douglas生産関数}で表されます。 \[
Y_t = \exp(a_t) \bar{K}^\alpha N_t^{1-\alpha} = \exp(a_t) N_t^{1-\alpha}, \qquad 0 \leq \alpha < 1.
\] この特定化のもとで、消費の限界効用 \(U_{C,t}\)、労働の限界効用
\(U_{N,t}\)、主観的限界代替率 \(MRS_t\)、および企業の労働限界生産力
\(MPN_t\) は、それぞれ以下のように明示的に計算されます。 \[
\begin{aligned}
U_{C,t} &= C_t^{-\gamma},\\
U_{N,t} &= -N_t^\varphi,\\
MRS_t &\equiv -\frac{U_{N,t}}{U_{C,t}} = N_t^\varphi C_t^\gamma,\\
MPN_t &\equiv (1-\alpha) \exp(a_t) N_t^{-\alpha}.
\end{aligned}
\] また、消費の限界効用に基づく確率的割引因子（SDF）は次式で表されます。
\[
Q_{t,t+1} = \beta \left( \frac{C_{t+1}}{C_t} \right)^{-\gamma}.
\]

\subsection{均衡労働投入量の閉形式解}\label{ux5747ux8861ux52b4ux50cdux6295ux5165ux91cfux306eux9589ux5f62ux5f0fux89e3}

労働市場の均衡条件
\(MRS_t = MPN_t\)（式~\ref{eq-labor-market}）に具体的な関数形を代入すると、次式が得られます。
\[
N_t^\varphi C_t^\gamma = (1-\alpha) \exp(a_t) N_t^{-\alpha}.
\] ここに財市場の清算条件 \(C_t = Y_t = \exp(a_t) N_t^{1-\alpha}\)
を代入して消費 \(C_t\) を消去します。 \[
N_t^\varphi \left[ \exp(a_t) N_t^{1-\alpha} \right]^\gamma = (1-\alpha) \exp(a_t) N_t^{-\alpha}.
\] 左辺を展開して整理すると、 \[
N_t^{\varphi + \gamma(1-\alpha)} \exp(\gamma a_t) = (1-\alpha) \exp(a_t) N_t^{-\alpha}
\] となり、\(N_t\) の項を左辺に、技術ショック \(a_t\)
の項を右辺にまとめると、 \[
N_t^{\varphi + \alpha + \gamma(1-\alpha)} = (1-\alpha) \exp\left[ (1-\gamma) a_t \right]
\] が導かれます。両辺を累乗根で解くことにより、外生的な技術水準 \(a_t\)
に対する均衡労働投入量 \(N_t\) の\textbf{閉形式解（Closed-form
Solution）}が一意に求まります。
\begin{equation}\protect\phantomsection\label{eq-labor-level}{
N_t = \left[ (1-\alpha) \exp\left\{ (1-\gamma) a_t \right\} \right]^{\frac{1}{\varphi + \alpha + \gamma(1-\alpha)}}.
}\end{equation} この解析解から明らかなように、技術ショック \(a_t\)
が労働投入量 \(N_t\) を増大させるか減少させるかは、指数部分に現れる
\((1-\gamma)\) の符号、すなわち相対的リスク回避度 \(\gamma\) が \(1\)
より大きいか小さいかによって一意に決定されます。

\subsection{構造パラメータの経済学的解釈}\label{ux69cbux9020ux30d1ux30e9ux30e1ux30fcux30bfux306eux7d4cux6e08ux5b66ux7684ux89e3ux91c8}

本モデルの挙動を規定する3つの主要な基礎パラメータ（\(\gamma, \varphi, \alpha\)）の経済学的意味は以下の通りです。

第一に、\(\gamma\) は\textbf{相対的リスク回避度（Relative Risk
Aversion）}であり、同時に\textbf{異時点間消費代替弾力性（EIS: Elasticity
of Intertemporal
Substitution）}の逆数（\(1/\gamma\)）を表します。確実性の世界において、債券オイラー方程式から
\(Q_{t,t+1} = 1/R_t\)
が成り立つため、実質利子率の変化に対する消費成長率の弾力性は \[
\frac{d\ln(C_{t+1}/C_t)}{d\ln R_t} = \frac{1}{\gamma}
\] となります。\(\gamma\)
が大きいほど（EISが小さいほど）、家計は異なる時点間での消費の偏りを嫌い、消費経路を滑らかに平準化しようとする選好が強くなります。

第二に、\(\varphi\)
は\textbf{Frisch（フリッシュ）労働供給弾力性}の逆数（\(1/\varphi\)）です。消費の限界効用
\(U_{C,t} = C_t^{-\gamma}\)
を一定に保持したもとで、実質賃金の上昇に対する労働供給の感応度を求めると、労働供給条件
\(W_t = N_t^\varphi C_t^\gamma\) から直ちに \[
\frac{d\ln N_t}{d\ln W_t}\bigg|_{U_C = \text{const}} = \frac{1}{\varphi}
\] が得られます。\(\varphi\)
が大きいほど（Frisch弾力性が小さいほど）、賃金が大きく変動しても労働供給量はあまり反応しなくなります。

第三に、\(\alpha\) は\textbf{資本分配率（Capital
Share）}を表します。本モデルでは資本ストックが \(\bar{K}=1\)
に固定されているため、\(\alpha\)
は生産された総付加価値のうち、固定資本の所有者（株主）に実質配当として帰属する割合（\(D_t = \alpha Y_t\)）を決定します。残りの割合
\((1-\alpha)\)
は労働者への実質賃金支払い（\(W_t N_t = (1-\alpha) Y_t\)）に充当されます。

\section{定常状態}\label{ux5b9aux5e38ux72b6ux614b}

対数線形近似を行うための基準点として、確率的ショックが存在しない確定的な「\textbf{定常状態（Steady
State）}」を求めます。外生的な技術ショックの定常水準を
\(a = 0\)（すなわち
\(A = \exp(0) = 1\)）と正規化します。定常状態では、すべての変数が時間を通じて一定の値（\(C_t = C, N_t = N, Y_t = Y, W_t = W, D_t = D\)）をとります。

\subsection{水準値における定常状態の導出}\label{ux6c34ux6e96ux5024ux306bux304aux3051ux308bux5b9aux5e38ux72b6ux614bux306eux5c0eux51fa}

定常状態における資源制約式、生産関数、企業の労働需要条件、および家計の労働供給条件は、それぞれ以下の通りです。
\[
\begin{aligned}
C &= Y = N^{1-\alpha},\\
W &= (1-\alpha) N^{-\alpha},\\
W &= N^\varphi C^\gamma.
\end{aligned}
\] 労働需要と労働供給を等置して実質賃金 \(W\) を消去し、さらに
\(C = N^{1-\alpha}\) を代入すると、 \[
N^\varphi \left( N^{1-\alpha} \right)^\gamma = (1-\alpha) N^{-\alpha} \implies N^{\varphi + \alpha + \gamma(1-\alpha)} = 1 - \alpha
\]
が得られます。したがって、定常状態における各変数の水準値は以下のように一意に求まります。
\[
\begin{aligned}
N &= (1-\alpha)^{\frac{1}{\varphi + \alpha + \gamma(1-\alpha)}},\\
Y &= C = N^{1-\alpha} = (1-\alpha)^{\frac{1-\alpha}{\varphi + \alpha + \gamma(1-\alpha)}},\\
W &= (1-\alpha) N^{-\alpha} = (1-\alpha)^{\frac{\varphi + \gamma(1-\alpha)}{\varphi + \alpha + \gamma(1-\alpha)}},\\
D &= \alpha Y = \alpha (1-\alpha)^{\frac{1-\alpha}{\varphi + \alpha + \gamma(1-\alpha)}}.
\end{aligned}
\] 配当 \(D = \alpha Y\) は固定資本 \(\bar{K}=1\)
に帰属する定常的な純収益を表しており、資本分配率 \(\alpha > 0\)
である限り常に正の値をとります。

ここで注意すべき点は、固定資本を \(1\)
に正規化したからといって、一般に定常状態の変数値が \(N = Y = C = W = 1\)
に揃うわけではないということです。資本分配率 \(\alpha > 0\)
の場合、\(1-\alpha < 1\) であるため、定常状態の労働投入量 \(N\) は \(1\)
未満の値となります。

唯一の例外は、資本シェアが存在しない\textbf{線形生産関数（\(\alpha = 0\)）}のケースです。このとき
\(Y_t = \exp(a_t) N_t\) かつ \(D_t = 0\)
となり、定常状態（\(a=0\)）において \[
N = 1, \qquad Y = C = 1, \qquad W = 1
\] が厳密に成立します。「定常状態の値がすべて \(1\)
になる」という直観は、この極めて特殊な線形生産のケースに限られた性質であることに留意してください。

\subsection{金融市場の定常状態}\label{ux91d1ux878dux5e02ux5834ux306eux5b9aux5e38ux72b6ux614b}

定常状態では消費水準が一定（\(C_{t+1} = C_t = C\)）であるため、家計の限界効用に基づく確率的割引因子の定常値は単に主観的割引因子と等しくなります。
\[
Q = \beta.
\] したがって、無リスク実質債券に関する消費オイラー方程式
\(1 = \mathbb{E}[Q R]\)
より、定常状態の実質総利子率は主観的時間選好率の逆数として定まります。
\[
R = \frac{1}{\beta}.
\] 中央銀行のインフレ目標や定常状態のインフレ率を
\(\Pi = 1\)（インフレ率ゼロ）と置けば、Fisher方程式を通じて、定常状態の名目総利子率も同様に
\(R^N = 1/\beta\) となります。

本講義ノートで以降使用する小文字の変数（\(a_t, y_t, c_t, n_t, w_t, r_t\)）は、すべてこれらの定常状態からの\textbf{対数偏差（パーセント乖離）}を表します。したがって、定常状態においてこれらはすべて
\(0\) となります。

\section{対数線形解と技術ショックの動学分析}\label{ux5bfeux6570ux7ddaux5f62ux89e3ux3068ux6280ux8853ux30b7ux30e7ux30c3ux30afux306eux52d5ux5b66ux5206ux6790}

\subsection{静学的均衡体系の対数線形化}\label{ux9759ux5b66ux7684ux5747ux8861ux4f53ux7cfbux306eux5bfeux6570ux7ddaux5f62ux5316}

定常状態の近傍で各均衡条件を1次のTaylor展開によって対数線形近似します。静学的均衡体系（式~\ref{eq-static-equilibrium}）を対数線形化すると、以下の4本の美しい線形連立方程式系が得られます。
\begin{equation}\protect\phantomsection\label{eq-loglinear-system}{
\begin{aligned}
y_t &= a_t + (1-\alpha) n_t,\\
c_t &= y_t,\\
w_t &= a_t - \alpha n_t,\\
w_t &= \gamma c_t + \varphi n_t.
\end{aligned}
}\end{equation}

これら4式は、それぞれ直観的な経済学的解釈を持ちます。

\begin{enumerate}
\def\labelenumi{\arabic{enumi}.}
\tightlist
\item
  \textbf{生産関数}（第1式）：産出の対数偏差 \(y_t\)
  は、全要素生産性の向上 \(a_t\) と、労働投入の増加 \((1-\alpha)n_t\)
  の和として決定されます。
\item
  \textbf{財市場清算}（第2式）：投資や政府支出が存在しないため、産出のパーセント変動はそのまま消費のパーセント変動と一致します（\(c_t = y_t\)）。
\item
  \textbf{労働需要}（第3式）：技術が \(1\%\)
  上昇すると、労働投入量を一定に保ったもとで企業の労働限界生産力（MPN）は
  \(1\%\)
  上昇します。一方、収穫逓減（\(\alpha > 0\)）のため、労働投入量が
  \(1\%\) 増加するとMPNは \(\alpha\%\) 低下します。
\item
  \textbf{労働供給}（第4式）：パラメータ \(\varphi\)
  は労働供給曲線の傾きを規定します。また、消費が \(1\%\)
  増加すると消費の限界効用が低下するため、家計は同じ労働を提供するためにより高い実質賃金を要求します。係数
  \(\gamma\)
  は、この消費増加に伴う「\textbf{所得効果}」の強さを表しています。
\end{enumerate}

\subsection{対数線形解の導出}\label{ux5bfeux6570ux7ddaux5f62ux89e3ux306eux5c0eux51fa}

連立方程式系（式~\ref{eq-loglinear-system}）を解くことで、技術ショック
\(a_t\) に対する各マクロ変数の反応を一意に導出できます。

まず、労働需要と労働供給の2式から実質賃金 \(w_t\) を消去すると、 \[
a_t - \alpha n_t = \gamma c_t + \varphi n_t
\] となります。ここに資源制約 \(c_t = y_t = a_t + (1-\alpha)n_t\)
を代入して消費 \(c_t\) を消去すると、 \[
a_t - \alpha n_t = \gamma \left[ a_t + (1-\alpha)n_t \right] + \varphi n_t
\] が得られます。\(n_t\) の項を左辺に、\(a_t\) の項を右辺に集めると、 \[
\left[ \varphi + \alpha + \gamma(1-\alpha) \right] n_t = (1-\gamma) a_t
\] が導かれます。したがって、均衡労働投入量の対数偏差 \(n_t\)
は次のように求まります。 \[
n_t = \frac{1-\gamma}{\varphi + \alpha + \gamma(1-\alpha)} a_t.
\]

ここで、産出および消費の技術ショックに対する反応弾力性を表す係数
\(\zeta_a\) を \[
\zeta_a \equiv \frac{1+\varphi}{\varphi + \alpha + \gamma(1-\alpha)} > 0
\] と定義します。このとき、産出 \(y_t\) は生産関数
\(y_t = a_t + (1-\alpha)n_t\) より、 \[
y_t = a_t + (1-\alpha) \frac{1-\gamma}{\varphi + \alpha + \gamma(1-\alpha)} a_t = \frac{\varphi + \alpha + \gamma(1-\alpha) + (1-\alpha)(1-\gamma)}{\varphi + \alpha + \gamma(1-\alpha)} a_t = \zeta_a a_t
\] と極めて簡潔に表されます。

同様に、実質賃金 \(w_t\) は労働需要式 \(w_t = a_t - \alpha n_t\) より、
\[
w_t = a_t - \alpha \frac{1-\gamma}{\varphi + \alpha + \gamma(1-\alpha)} a_t = \frac{\varphi + \gamma}{\varphi + \alpha + \gamma(1-\alpha)} a_t = \frac{1 - \alpha \zeta_a}{1-\alpha} a_t
\] として導出されます。

以上をまとめると、対数線形化された均衡解は以下の明快な体系として得られます。
\begin{equation}\protect\phantomsection\label{eq-loglinear-solution}{
\begin{aligned}
y_t &= \zeta_a a_t,\\
c_t &= y_t,\\
n_t &= \frac{1-\gamma}{\varphi + \alpha + \gamma(1-\alpha)} a_t = \frac{\zeta_a - 1}{1-\alpha} a_t,\\
w_t &= \frac{\varphi + \gamma}{\varphi + \alpha + \gamma(1-\alpha)} a_t = \frac{1-\alpha\zeta_a}{1-\alpha} a_t.
\end{aligned}
}\end{equation} 分母の \(\varphi + \alpha + \gamma(1-\alpha)\)
は、労働供給の傾き（\(\varphi\)）、企業の収穫逓減度（\(\alpha\)）、および消費の限界効用低下を通じた所得効果（\(\gamma(1-\alpha)\)）の3要素が合算されたものであり、経済の安定性を保証する正の定数です。

\subsection{技術ショックの動学とインパルス応答}\label{ux6280ux8853ux30b7ux30e7ux30c3ux30afux306eux52d5ux5b66ux3068ux30a4ux30f3ux30d1ux30ebux30b9ux5fdcux7b54}

外生的な技術ショック \(a_t\)
は、以下の標準的な1階自己回帰過程（AR(1)過程）に従って推移すると仮定します。
\begin{equation}\protect\phantomsection\label{eq-technology-process}{
a_t = \rho_a a_{t-1} + \epsilon_t^a, \qquad 0 \leq \rho_a < 1, \qquad \mathbb{E}_{t-1}[\epsilon_t^a] = 0.
}\end{equation} ここで \(\rho_a\)
は技術ショックの\textbf{持続性（Persistence）}を表すパラメータであり、\(\epsilon_t^a\)
は平均ゼロ、分散 \(\sigma_a^2\) の白色雑音（White
Noise）イノベーションです。

各期の期首時点において、前末までに持ち越された状態変数 \(a_{t-1}\)
が所与として存在し、そこに新たなサプライズ・ショック \(\epsilon_t^a\)
が加わることで、当期の技術水準 \(a_t\)
が確定・観察されます。家計と企業は、この \(a_t\)
を含む当期の情報集合に基づいて最適な消費・労働・生産を決定します。

時点 \(t\) において経済が定常状態（\(a_{t-1}=0\)）にあるとき、大きさ
\(1\%\)（\(\epsilon_t^a = 0.01\)）の技術ショックが発生し、その後のショックがゼロ（\(\epsilon_{t+h}^a = 0, h \geq 1\)）であったとします。このとき、\(h\)
期間先における技術水準および各内生変数の\textbf{インパルス応答（Impulse
Response Function:
IRF）}は、以下のように指数関数的に減衰する経路を描きます。
\begin{equation}\protect\phantomsection\label{eq-technology-irf}{
a_{t+h} = \rho_a^h \epsilon_t^a, \qquad x_{t+h} = \kappa_x \rho_a^h \epsilon_t^a \quad (x \in \{y, c, n, w\}).
}\end{equation} ここで各変数の反応係数 \(\kappa_x\)
は、対数線形解（式~\ref{eq-loglinear-solution}）から直ちに求まります。すなわち、
\[
\kappa_y = \zeta_a, \qquad \kappa_c = \zeta_a, \qquad \kappa_n = \frac{1-\gamma}{\varphi+\alpha+\gamma(1-\alpha)}, \qquad \kappa_w = \frac{\varphi+\gamma}{\varphi+\alpha+\gamma(1-\alpha)}
\]
となります。すべてのマクロ変数は、ショック発生期（\(h=0\)）に最も大きく反応し、その後は共通の減衰率
\(\rho_a\) に従って滑らかに定常状態へと復帰します。

\subsection{代替効果と所得効果の対立}\label{ux4ee3ux66ffux52b9ux679cux3068ux6240ux5f97ux52b9ux679cux306eux5bfeux7acb}

対数線形解から得られる最も深い理論的洞察は、技術進歩が労働投入量 \(n_t\)
に与える影響の方向性が、家計の選好パラメータ \(\gamma\)
の値によって決定されるという点です。

技術水準 \(a_t\)
が向上すると、労働の限界生産力（MPN）が高まり、市場の実質賃金 \(w_t\)
が上昇します。この賃金上昇は、家計に対して以下の2つの相反する力を及ぼします。

\begin{enumerate}
\def\labelenumi{\arabic{enumi}.}
\tightlist
\item
  \textbf{代替効果（Substitution Effect）}：
  実質賃金の上昇は、余暇を過ごす機会費用の上昇を意味します。したがって、余暇を減らして労働供給を増やそうとする誘因が働きます（労働投入に対する正の効果）。
\item
  \textbf{所得効果（Income Effect）}：
  賃金上昇および企業からの配当収入の増加は、家計の実質生涯富（購買力）を増大させます。余暇が正常財であるならば、豊かになった家計はより多くの余暇を消費しようとするため、労働供給を減らそうとする誘因が働きます（労働投入に対する負の効果）。
\end{enumerate}

この2つの効果の力関係は、相対的リスク回避度（EISの逆数）\(\gamma\)
によって完全に決まります。

{\def\LTcaptype{none} % do not increment counter
\begin{longtable}[]{@{}
  >{\raggedright\arraybackslash}p{(\linewidth - 8\tabcolsep) * \real{0.2000}}
  >{\raggedright\arraybackslash}p{(\linewidth - 8\tabcolsep) * \real{0.2000}}
  >{\raggedright\arraybackslash}p{(\linewidth - 8\tabcolsep) * \real{0.2000}}
  >{\raggedright\arraybackslash}p{(\linewidth - 8\tabcolsep) * \real{0.2000}}
  >{\raggedright\arraybackslash}p{(\linewidth - 8\tabcolsep) * \real{0.2000}}@{}}
\toprule\noalign{}
\begin{minipage}[b]{\linewidth}\raggedright
\(\gamma\) の領域
\end{minipage} & \begin{minipage}[b]{\linewidth}\raggedright
支配的な効果
\end{minipage} & \begin{minipage}[b]{\linewidth}\raggedright
労働投入 \(n_t\) の反応
\end{minipage} & \begin{minipage}[b]{\linewidth}\raggedright
産出 \(y_t\) の反応係数 \(\zeta_a\)
\end{minipage} & \begin{minipage}[b]{\linewidth}\raggedright
経済学的直観
\end{minipage} \\
\midrule\noalign{}
\endhead
\bottomrule\noalign{}
\endlastfoot
\textbf{\(\gamma < 1\)} & \textbf{代替効果が支配的} &
\textbf{増加（\(n_t > 0\)）} & \textbf{\(\zeta_a > 1\)} &
余暇の異時点間代替を積極的に行うため、賃金上昇に強く反応して就業時間を伸ばす。労働増加により産出は技術ショック以上に拡大する。 \\
\textbf{\(\gamma = 1\)} & \textbf{完全に相殺} &
\textbf{不変（\(n_t = 0\)）} & \textbf{\(\zeta_a = 1\)} &
対数効用（\(\ln C\)）のもとでは代替効果と所得効果が完璧に打ち消し合い、労働時間は技術ショックに対して全く動かない。 \\
\textbf{\(\gamma > 1\)} & \textbf{所得効果が支配的} &
\textbf{減少（\(n_t < 0\)）} & \textbf{\(\zeta_a < 1\)} &
消費平準化の選好が強く、富の増加を余暇の拡大（労働の削減）へと振り向ける。労働減少のため産出の伸びは技術向上率を下回る。 \\
\end{longtable}
}

具体的な数値例として、標準的なパラメータ値（\(\alpha = 0.4, \varphi = 1.0\)）のもとで、\(\gamma\)
の違いによるインパクト反応係数を比較すると、以下のようになります。

{\def\LTcaptype{none} % do not increment counter
\begin{longtable}[]{@{}
  >{\raggedleft\arraybackslash}p{(\linewidth - 6\tabcolsep) * \real{0.2500}}
  >{\raggedleft\arraybackslash}p{(\linewidth - 6\tabcolsep) * \real{0.2500}}
  >{\raggedleft\arraybackslash}p{(\linewidth - 6\tabcolsep) * \real{0.2500}}
  >{\raggedleft\arraybackslash}p{(\linewidth - 6\tabcolsep) * \real{0.2500}}@{}}
\toprule\noalign{}
\begin{minipage}[b]{\linewidth}\raggedleft
相対的リスク回避度 \(\gamma\)
\end{minipage} & \begin{minipage}[b]{\linewidth}\raggedleft
産出反応係数 \(\kappa_y = \zeta_a\)
\end{minipage} & \begin{minipage}[b]{\linewidth}\raggedleft
労働反応係数 \(\kappa_n\)
\end{minipage} & \begin{minipage}[b]{\linewidth}\raggedleft
賃金反応係数 \(\kappa_w\)
\end{minipage} \\
\midrule\noalign{}
\endhead
\bottomrule\noalign{}
\endlastfoot
\textbf{\(\gamma = 0.5\)} & \(1.176\) & \(+0.294\) & \(0.882\) \\
\textbf{\(\gamma = 1.0\)} & \(1.000\) & \(0.000\) & \(1.000\) \\
\textbf{\(\gamma = 2.0\)} & \(0.769\) & \(-0.385\) & \(1.154\) \\
\end{longtable}
}

\begin{tcolorbox}[enhanced jigsaw, arc=.35mm, bottomrule=.15mm, bottomtitle=1mm, breakable, colback=white, colbacktitle=quarto-callout-important-color!10!white, colframe=quarto-callout-important-color-frame, coltitle=black, left=2mm, leftrule=.75mm, opacityback=0, opacitybacktitle=0.6, rightrule=.15mm, title=\textcolor{quarto-callout-important-color}{\faExclamation}\hspace{0.5em}{資本蓄積のないモデル特有の性質に関する注意}, titlerule=0mm, toprule=.15mm, toptitle=1mm]

労働の反応が \((1-\gamma)\)
の符号のみで決まり、後述のように正の一時的技術ショックに対して自然実質利子率が低下するという結果は、「\textbf{投資による資本蓄積が存在しない}」「\textbf{消費と労働の効用が加法分離可能である}」「\textbf{財市場清算が常に
\(C_t = Y_t\) である}」という本モデル固有の単純化に強く依存しています。

次学期以降に学ぶ資本蓄積を含む標準的なRBCモデルでは、技術向上に伴って投資需要が大きく喚起されるため、たとえ
\(\gamma > 1\)
であっても労働投入や実質金利が一時的に上昇するなど、より複雑で豊かな動学が現れます。

\end{tcolorbox}

\section{自然実質利子率と古典派の二分法}\label{ux81eaux7136ux5b9fux8ceaux5229ux5b50ux7387ux3068ux53e4ux5178ux6d3eux306eux4e8cux5206ux6cd5}

\subsection{実物配分の先行決定と金融市場}\label{ux5b9fux7269ux914dux5206ux306eux5148ux884cux6c7aux5b9aux3068ux91d1ux878dux5e02ux5834}

前章で確認した通り、実物配分（\(n_t, y_t, c_t, w_t\)）は、家計の労働供給、企業の労働需要、生産技術、および総資源制約の4条件のみから完全に先行して決定されます。

民間債券は家計間で貸借される純供給ゼロの資産であるため、金融市場において決定される実質総利子率
\(R_t\) は、実物側で既に決まった消費成長経路（\(C_t\) から \(C_{t+1}\)
への変化）と、債券市場の需給清算（\(B_t = 0\)）とを整合させるように受動的に決まる「市場清算価格」としての役割を果たします。

\subsection{消費オイラー方程式と自然実質利子率の導出}\label{ux6d88ux8cbbux30aaux30a4ux30e9ux30fcux65b9ux7a0bux5f0fux3068ux81eaux7136ux5b9fux8ceaux5229ux5b50ux7387ux306eux5c0eux51fa}

家計の無リスク実質債券に関する消費オイラー方程式は、次式で与えられます。
\begin{equation}\protect\phantomsection\label{eq-real-bond-euler}{
1 = \mathbb{E}_t \left[ \beta \left( \frac{C_{t+1}}{C_t} \right)^{-\gamma} R_t \right].
}\end{equation} \(R_t\) は時点 \(t\)
において確定する1期間の無リスク実質総利子率であるため、期待値の外に取り出すことができ、
\[
R_t = \left\{ \mathbb{E}_t \left[ \beta \left( \frac{C_{t+1}}{C_t} \right)^{-\gamma} \right] \right\}^{-1}
\] と表現されます。

この式を定常状態（\(R = 1/\beta, C_{t+1}=C_t=C\)）の周りで1次対数近似します。実質総利子率の対数偏差を
\(r_t \equiv \ln(R_t / R) \approx R_t - R\)、主観的時間選好率を
\(\rho \equiv -\ln\beta \approx 1/\beta - 1\)
と定義すると、以下の対数線形消費オイラー方程式が得られます。
\begin{equation}\protect\phantomsection\label{eq-natural-rate-euler}{
r_t = \gamma \mathbb{E}_t [c_{t+1} - c_t].
}\end{equation} この方程式は、経済全体の\textbf{期待消費成長率
\(\mathbb{E}_t[c_{t+1} - c_t]\)}
が高いときほど、家計が当期の消費を我慢して貯蓄を行うために、より高い実質利子率
\(r_t\) が市場で要求されることを意味しています。また、相対的リスク回避度
\(\gamma\) が大きいほど（EIS \(1/\gamma\)
が小さいほど）、消費を平準化しようとする動機が強くなるため、同じ期待消費成長率を支えるために必要な利子率の変動幅は大きくなります。

ここで、外生ショック過程（式~\ref{eq-technology-process}）より
\(\mathbb{E}_t[a_{t+1}] = \rho_a a_t\) であり、また均衡消費が
\(c_t = \zeta_a a_t\) であることから、翌期の期待消費成長率は \[
\mathbb{E}_t [c_{t+1} - c_t] = \zeta_a \mathbb{E}_t [a_{t+1} - a_t] = \zeta_a (\rho_a - 1) a_t
\]
と計算されます。これをオイラー方程式（式~\ref{eq-natural-rate-euler}）に代入すると、価格が完全に伸縮的な経済において実現する実質利子率である「\textbf{自然実質利子率（Natural
Rate of Interest: \(r_t^n\)）}」の解析解が得られます。
\begin{equation}\protect\phantomsection\label{eq-natural-rate-technology}{
r_t^n = \gamma \zeta_a (\rho_a - 1) a_t.
}\end{equation}

ショックの持続性が \(0 \leq \rho_a < 1\)
であるとき、\((\rho_a - 1) < 0\)
となります。したがって、\textbf{正の一時的技術ショック（\(a_t > 0\)）が発生すると、自然実質利子率は必ず低下（\(r_t^n < 0\)）する}という重要な結果が導かれます。

この現象の直観は以下の通りです。正の技術ショックが発生した当期は生産性が高く、消費水準
\(c_t\)
が一時的に急上昇します。しかし、ショックが将来に向けて徐々に減衰していくため、家計は「当期の消費水準に比べて、将来の消費水準は相対的に低くなる（負の消費成長率：\(\mathbb{E}_t[\Delta c_{t+1}] < 0\)）」と予想します。強い消費平準化動機を持つ家計は、恵まれた当期の所得を将来のために貯蓄しようと殺到しますが、経済全体では貯蓄手段（資本蓄積）が存在せず、債券の純供給もゼロに固定されています。その結果、家計の過剰な貯蓄意欲を抑制し、当期のうちにすべての生産物を消費し尽くさせるために、市場の実質利子率が大幅に低下しなければならないのです。

ショックが恒久的（\(\rho_a \to 1\)）に近づくほど、当期の豊かさと将来の豊かさが等しくなるため、消費成長率の低下予想は弱まり、自然実質利子率の低下幅はゼロへと近づきます。

\subsection{古典派の二分法とお金の中立性}\label{ux53e4ux5178ux6d3eux306eux4e8cux5206ux6cd5ux3068ux304aux91d1ux306eux4e2dux7acbux6027}

実物サイドの方程式群は、実質配分（\(n_t, y_t, c_t, w_t\)）および実質利子率
\(r_t^n\)
を完全に一意に決定します。しかし、名目変数である\textbf{名目総利子率
\(R_t^N\)}、\textbf{粗インフレ率
\(\Pi_{t+1} \equiv P_{t+1}/P_t\)}、および\textbf{絶対物価水準 \(P_t\)}
は、この実物ブロックの中からは一切決定されません。

家計の名目債券保有に関するオイラー方程式は、以下のように記述されます。
\begin{equation}\protect\phantomsection\label{eq-nominal-bond-euler}{
1 = \mathbb{E}_t \left[ Q_{t,t+1} \frac{R_t^N}{\Pi_{t+1}} \right].
}\end{equation} ここで \(R_t^N / \Pi_{t+1}\)
は、翌期の実現インフレ率によって実質化された事後的な実質収益率を表します。

実質配分が名目変数（貨幣供給量や名目価格水準）の決定から完全に独立しているというこの性質を、マクロ経済学では「\textbf{古典派の二分法（Classical
Dichotomy）}」と呼びます。

実物配分と名目体系を結びつけ、具体的なインフレ率や物価水準を一意に決定するためには、中央銀行の金融政策ルール（Taylorルールなど）や貨幣市場の均衡条件をモデルに追加する必要があります。
- \textbf{価格が完全に伸縮的な場合}： 物価水準 \(P_t\) やインフレ率
\(\Pi_t\)
が瞬時に調整されるため、中央銀行がどのような政策をとろうとも、実際の実質利子率は常に自然実質利子率と一致（\(r_t = r_t^n\)）し、実物配分は不変に保たれます（貨幣の中立性）。
- \textbf{価格に硬直性が存在する場合（New
Keynesianモデル：第7回以降）}：
企業の価格改定が遅れるため、古典派の二分法は崩壊します。中央銀行が名目利子率を操作すると、実際の実質利子率
\(r_t\) が自然実質利子率 \(r_t^n\)
から乖離し、その差（実質金利ギャップ）が産出量と自然産出量の乖離である「\textbf{需給ギャップ（産出ギャップ）}」を生み出すことになります。

\section{まとめと第5回への接続}\label{ux307eux3068ux3081ux3068ux7b2c5ux56deux3078ux306eux63a5ux7d9a}

本講義では、固定資本を \(1\)
に正規化し、投資による資本蓄積を行わない最も純粋な柔軟価格RBCモデルを構築・分析しました。本講義の要点は以下の通りです。

\begin{enumerate}
\def\labelenumi{\arabic{enumi}.}
\tightlist
\item
  \textbf{実物配分の一意性}：
  実物配分の核心は、労働市場における主観的限界代替率と技術的限界生産力の一致条件
  \(MRS_t = W_t = MPN_t\)（式~\ref{eq-labor-market}）にあります。
\item
  \textbf{代替効果と所得効果}：
  技術ショックに対する産出、消費、実質賃金の反応は常に正ですが、労働投入量の反応方向は相対的リスク回避度
  \(\gamma \gtrless 1\)（代替効果と所得効果の大小関係）によって決定されます。
\item
  \textbf{自然実質利子率の受動的決定}： 自然実質利子率 \(r_t^n\)
  は、実物側で先に決まった消費成長経路と整合するように消費オイラー方程式から受動的に復元されます。一時的な正の技術ショックは、消費平準化の動機を通じて自然実質利子率を低下させます。
\end{enumerate}

\subsection{第5回への接続：政府支出の導入}\label{ux7b2c5ux56deux3078ux306eux63a5ux7d9aux653fux5e9cux652fux51faux306eux5c0eux5165}

次回の第5回講義では、本モデルに\textbf{政府部門（政府支出 \(G_t\)
および一括固定税 \(T_t\)）}を導入します。

総資源制約式は、本講義の \(Y_t = C_t\) から \[
Y_t = C_t + G_t
\]
へと拡張されます。固定資本が蓄積できないという前提は維持されるため、政府支出の拡大は民間投資をクラウドアウトするのではなく、同一期内における家計の消費と労働供給の調整（負の富効果を通じた労働供給の増加）によって賄われます。

また、本講義の債券 \(B_t\)
は家計間で貸借される純供給ゼロの民間債券でしたが、第5回では政府が発行する国債が登場し、金融資産が正の純供給を持つ点に構造的な違いがあります。なお、第5回の一般化されたモデルにおいて定常状態の政府支出比率をゼロ（\(G/Y = 0\)）とし、政府支出ショックをゼロと置けば、本講義で導出したすべての解体系に厳密に一致します。

\section{参考文献}\label{ux53c2ux8003ux6587ux732e}

\textbf{基本文献}

\begin{itemize}
\tightlist
\item
  Galí, Jordi (2015) \emph{Monetary Policy, Inflation, and the Business
  Cycle: An Introduction to the New Keynesian Framework and Its
  Applications} (2nd ed.). Princeton University Press.
  \textbf{第2章が資本蓄積のない柔軟価格モデルと自然実質利子率の最も標準的な参照先}
\item
  Romer, David (2019) \emph{Advanced Macroeconomics} (5th ed.).
  McGraw-Hill.
  \textbf{実物ショックの伝播メカニズムと標準的な動学モデルの比較}
\end{itemize}

\textbf{発展文献：資本蓄積を含む古典的RBCモデル}

\begin{itemize}
\tightlist
\item
  Kydland, Finn E., and Edward C. Prescott (1982) ``Time to Build and
  Aggregate Fluctuations.'' \emph{Econometrica}, 50(6), pp.~1345-1370.
\item
  Hansen, Gary D. (1985) ``Indivisible Labor and the Business Cycle.''
  \emph{Journal of Monetary Economics}, 16(3), pp.~309-327.
\item
  King, Robert G., Charles I. Plosser, and Sergio T. Rebelo (1988)
  ``Production, Growth and Business Cycles: I. The Basic Neoclassical
  Model.'' \emph{Journal of Monetary Economics}, 21(2-3), pp.~195-232.
\item
  King, Robert G., and Sergio T. Rebelo (1999) ``Resuscitating Real
  Business Cycles.'' In John B. Taylor and Michael Woodford (eds.),
  \emph{Handbook of Macroeconomics}, Vol. 1B, ch.~14, pp.~927-1007.
  Elsevier.
\item
  Plosser, Charles I. (1989) ``Understanding Real Business Cycles.''
  \emph{Journal of Economic Perspectives}, 3(3), pp.~51-77.
\end{itemize}

\section{演習問題}\label{ux6f14ux7fd2ux554fux984c}

\subsection{基本問題}\label{ux57faux672cux554fux984c}

\textbf{問1
［Core］：均衡概念の直観的理解}　資本蓄積のないRBCモデルにおいて、労働市場の均衡条件
\(MRS_t = W_t = MPN_t\)
が成立することの経済学的意味を、家計の最適化行動と企業の利潤最大化行動の双方の観点から平易な言葉で説明しなさい。

\textbf{問2 ［Core］［Derivation］：閉形式解の導出}　効用関数
\(U(C_t, N_t) = \frac{C_t^{1-\gamma}-1}{1-\gamma} - \frac{N_t^{1+\varphi}}{1+\varphi}\)、生産関数
\(Y_t = \exp(a_t) \bar{K}^\alpha N_t^{1-\alpha}\)（\(\bar{K}=1\)）、および財市場清算条件
\(Y_t = C_t\)
から、均衡労働投入量の水準値を与える閉形式解（式~\ref{eq-labor-level}）を途中の計算過程を含めて厳密に導出しなさい。

\textbf{問3
［Core］：技術ショックの波及経路}　正の技術ショック（\(a_t > 0\)）が発生したとき、産出量、消費水準、および実質賃金がどのように変動するかを説明しなさい。また、労働投入量の反応の符号が相対的リスク回避度
\(\gamma\)
の値に依存する理由を、代替効果と所得効果の対立に触れて述べなさい。

\textbf{問4
［Core］［Derivation］：固定資本と企業利潤}　企業の労働需要条件とCobb-Douglas生産関数から、総賃金支払い額が
\(W_t N_t = (1-\alpha) Y_t\)、実質配当（利潤）が \(D_t = \alpha Y_t\)
となることを示しなさい。また、線形生産関数（\(\alpha = 0\)）の極限において、配当額および企業の株式価値がどのようになるかを論じなさい。

\subsection{発展問題}\label{ux767aux5c55ux554fux984c}

\textbf{問5 ［Research
extension］［Computation］：インパルス応答の具体的計算}　パラメータが
\(\alpha = 0.4, \varphi = 1.0, \gamma = 2.0, \rho_a = 0.9\)
である経済を考えます。定常状態にある当期（\(t\)）に大きさ
\(1\%\)（\(0.01\)）の正の技術ショックが発生し、その後のショックがゼロであったとします。このとき、技術水準
\(a_{t+h}\)、産出量 \(y_{t+h}\)、労働投入量 \(n_{t+h}\)、および実質賃金
\(w_{t+h}\) のインパルス応答の値を、\(h = 0, 1, 2\)
についてそれぞれ計算しなさい。

\textbf{問6 ［Research
extension］［Derivation］：自然実質利子率の導出と持続性}　対数線形消費オイラー方程式
\(r_t = \gamma \mathbb{E}_t [c_{t+1} - c_t]\) と均衡消費関数
\(c_t = \zeta_a a_t\)
から、自然実質利子率の決定式（式~\ref{eq-natural-rate-technology}）を導出しなさい。また、技術ショックの持続性パラメータ
\(\rho_a\) が \(1\)
に近づくにつれて、技術ショックに対する自然実質利子率の反応幅が小さくなる経済学的理由を直観的に説明しなさい。

\begin{samepage}

\textbf{問7 ［Research
extension］：第5回モデルへの拡張と構造比較}　政府支出を導入した総資源制約式
\(Y_t = C_t + G_t\)
において、定常状態の政府支出シェアをゼロ（\(G/Y = 0\)）、政府支出ショックをゼロ（\(g_t = 0\)）とした極限が、本講義の資源制約式
\(Y_t = C_t\)
と一致することを示しなさい。また、民間債券の純供給がゼロである本講義の金融市場と、正の純供給を持つ政府債務が存在する第5回の金融市場の違いが、家計の富やマクロ均衡にどのような相違をもたらすかについて論じなさい。

\end{samepage}




\end{document}
